\begin{helper}
Let \(V\) be finite and let \(\nu\) be a measure on activation--priority
pairs \((A,\pi)\in\mathcal P(V)\times\mathbb R^V\).  Fix
\(0\le p\le1\).  Assume that every activation atom satisfies
\[
\nu\{A=A_0\}=
\operatorname{ofReal}\!\left(p^{|A_0|}(1-p)^{|V|-|A_0|}\right)
\]
for every \(A_0\subseteq V\).  Assume also that activated priorities are
supported in the unit cube on each atom, in the sense that
\[
\nu\{A=A_0,\ \pi(q)\in[0,1]\text{ for every }q\in A_0\}
=\nu\{A=A_0\}
\]
for every \(A_0\subseteq V\).  Finally, assume the full lower-rectangle
product law: for every \(A_0\subseteq V\) and every threshold vector
\(t:V\to\mathbb R\) with \(t(q)\in[0,1]\) for all \(q\),
\[
\nu\{A=A_0,\ 0\le \pi(q)\le t(q)\text{ for every }q\in V\}
=\nu\{A=A_0\}\operatorname{ofReal}\!\left(\prod_{q\in V}t(q)\right).
\]
Let \(u\ne v\), and let \(U,T\subseteq V\) with
\(u\notin U\cup T\) and \(v\notin U\cup T\).  Then for every \(z\in U\cup T\),
\[
\nu\{u,v,z\in A,\ \pi(q)\in[0,1]\text{ for every }q\in A,\
\pi(z)=\pi(u)\text{ or }\pi(z)=\pi(v)\}=0 .
\]
\end{helper}
\begin{proof}
Take \(z\in U\cup T\).  Since \(u\notin U\cup T\) and \(v\notin U\cup T\), we
have \(z\ne u\) and \(z\ne v\).  We prove the required nullity by first
working on each activation atom.

Fix an activation atom \(A_0\), and put
\[
\alpha(A_0)=p^{|A_0|}(1-p)^{|V|-|A_0|}.
\]
The activation-atom hypothesis gives
\(\nu\{A=A_0\}=\operatorname{ofReal}(\alpha(A_0))\), with
\(\alpha(A_0)\ge0\) because \(0\le p\le1\).  Let
\[
Q=\{(A,\pi): A=A_0,\ 0\le \pi(q)\le1\hbox{ for every }q\in V\}.
\]
The lower-rectangle hypothesis with constant threshold \(1\) gives
\[
\nu(Q)=\nu\{A=A_0\}\operatorname{ofReal}\!\left(\prod_{q\in V}1\right)
=\nu\{A=A_0\},
\]
so the part of the atom outside the full cube \(Q\) has \(\nu\)-measure
\(0\).

On \(Q\), the lower-rectangle hypothesis determines the atomwise priority
law on all anchored rectangles.  More explicitly, for every threshold vector
\(t:V\to[0,1]\), the measure of the anchored rectangle
\(\prod_{q\in V}[0,t(q)]\) inside the atom is
\[
\operatorname{ofReal}(\alpha(A_0))\,
\operatorname{ofReal}\!\left(\prod_{q\in V}t(q)\right).
\]
This is the same value assigned by
\(\operatorname{ofReal}(\alpha(A_0))\) times finite product Lebesgue measure
on \([0,1]^V\).  The anchored rectangles are a finite-intersection-stable
generating class for the Borel sets of the finite cube, and both finite
measures have total mass \(\operatorname{ofReal}(\alpha(A_0))\).  The
finite-measure uniqueness theorem therefore identifies the atom-restricted
priority law on the full cube with
\(\operatorname{ofReal}(\alpha(A_0))\lambda_V\), where \(\lambda_V\) is
finite product Lebesgue measure on \([0,1]^V\).

If \(A_0\) does not contain all three of \(u,v,z\), then the atomwise part of
the boundary event is empty.  Suppose now that \(u,v,z\in A_0\).  Because
\(z\ne u\), the set
\[
D_u=\{\pi\in[0,1]^V:\pi(z)=\pi(u)\}
\]
is a coordinate hyperplane.  By Fubini over the \(z\)-coordinate, after all
other coordinates have been fixed the section of \(D_u\) is the singleton
\(\{\pi(u)\}\), which has one-dimensional Lebesgue measure \(0\).  This is
the fixed-section zero-length interval consequence of
\noderef{SamplingFinitePriorityOneCoordinateIntervalMass}; the node is not
being used as a variable-coordinate hyperplane theorem, only as the
one-dimensional interval calculation after the Fubini slice is fixed.  Hence
\(\lambda_V(D_u)=0\), and the atomwise full-cube event with
\(\pi(z)=\pi(u)\) has \(\nu\)-measure \(0\).

The same argument applies to
\[
D_v=\{\pi\in[0,1]^V:\pi(z)=\pi(v)\},
\]
because \(z\ne v\).  Therefore the atomwise full-cube event with
\(\pi(z)=\pi(v)\) also has \(\nu\)-measure \(0\).  The union of these two
atomwise full-cube boundary pieces has measure \(0\) by finite
subadditivity.

Finally compare the full cube with the activated-coordinate cube appearing
in the theorem.  The cube-support hypothesis gives
\[
\nu\{A=A_0,\ \pi(q)\in[0,1]\hbox{ for every }q\in A_0\}
=\nu\{A=A_0\},
\]
so, inside the atom, the complement of the activated-coordinate cube has
measure \(0\).  The atomwise boundary event from the statement is contained
in the union of this zero-measure complement with the two full-cube
hyperplane events just proved null.  Thus every activation atom contributes
measure \(0\).  Since \(V\) is finite, there are finitely many activation
atoms, and the whole boundary event for this fixed \(z\) is their finite
union.  Finite subadditivity gives the claimed total measure \(0\).
\end{proof}
